\section{Restoring a File}\label{sec:restore}

\subsection{Writing the file}
A restore starts only for a ghost in state \state{ready}, into a folder the user chooses, under a
name that may not exist there and must fit in 255 bytes. The record moves to \state{restoring} and
names the destination and a hidden temporary file beside it,
\code{.<name>.<12\nobreakspace{}hex>.ghostvault-part}, created exclusively and readable by the owner only.

Chunks are downloaded in parallel, eight at a time, with at most sixteen held in memory, and are
written into the temporary file in order. Each is checked as described in \secref{sec:checks}
before it is written, and the first failure stops every further download. After the last chunk the
whole file's SHA-256 is compared with the original's. The file then gets the original's permission
bits (without setuid, setgid or sticky bits) and its access and modification times, or on Windows
its read-only attribute and times, and is flushed to disk. Nothing has a final name yet; a failure
up to here removes the temporary file and leaves the ghost \state{ready}.

\subsection{Publishing without replacing}\label{sec:publish}
The finished file is given its final name in a way that can never replace a file, however recently
that file appeared. On Linux and macOS this is a hard link to the final name, which fails if the
name exists, followed by removal of the temporary name; on file systems without hard links it is
\code{renameat2} with \code{RENAME\_NOREPLACE} (Linux) or \code{renamex\_np} with
\code{RENAME\_EXCL} (macOS), and where neither exists the restore is refused rather than risk an
overwrite. On Windows it is \code{MoveFileExW} without the replace flag and with write-through. If a
file of that name appeared during the restore, it is left untouched, the temporary file is removed,
and the user is asked for another name. The folder is then flushed, and only then does the record
move to \state{cleanup}.

\subsection{Removing the remote copy}
Clean-up removes the ghost's paths from each repository with one commit, then reads the branch again;
the removal is complete only when a fresh listing shows none of them. A round that GitHub refused,
applied only in part, or acknowledged without effect is followed by another, after 2, 4 and 8\,s,
up to four rounds. Rounds can be repeated safely: the paths are random and recorded, absent ones are
skipped, and the branch is never forced. A refusal that repeating cannot fix (a rejected token, a
missing repository, an exhausted rate limit) ends the rounds at once. If paths are still listed after
the last round, the error says how many, with GitHub's last status, message and request id and what
each round sent and saw; the ghost stays in \state{cleanup}, is shown with a Retry button, and is
retried at the next start.

Only when every repository has confirmed the removal are the record and its cached preview deleted.
Before the record is deleted, its id is added to the ghosts GhostVault has removed (when a recovery
file or the repository backup may hold it), so that a later backup can leave it out without the
removal being mistaken for a lost record (\secref{sec:incomplete}).

Removal takes the ghost off the branch. It does not erase it: the commits that added the chunks
remain in the repository's history, and GitHub may keep unreachable objects for a time. A removed
ghost stays unreadable only as long as its key exists nowhere: not in the data folder, not in a
backup of that folder, and not in a recovery bundle made while the ghost existed.
